\documentclass[11pt]{article}

\usepackage[margin=1in]{geometry}
\usepackage{amsmath,amssymb,amsthm,mathtools}
\usepackage{microtype}

\newtheorem{theorem}{Theorem}
\newtheorem{lemma}{Lemma}
\newtheorem{corollary}{Corollary}
\theoremstyle{remark}
\newtheorem{remark}{Remark}

\newcommand{\E}{\mathbb{E}}
\renewcommand{\Pr}{\mathbb{P}}
\newcommand{\Ind}[1]{\mathbb{I}\big\{{#1}\big\}}
\newcommand{\bp}{\boldsymbol{p}}
\newcommand{\cV}{\mathcal{V}}

\title{A Logarithmic Lower Bound for Adversarial Delays}
\author{}
\date{}

\begin{document}
\maketitle

\begin{theorem}
\label{thm:log-lower-bound}
For every horizon $T\ge2$ and every possibly randomized online policy
$\pi$, there exists a deterministic, oblivious, nonincreasing price
sequence
\[
1=p_1\ge p_2\ge\cdots\ge p_T>0
\]
such that
\[
\cV_T=0
\]
and
\[
R_T
\ge
\frac{3}{16}\bigl(H_T-1\bigr).
\]
\end{theorem}

\begin{proof}

The proof uses a distribution over deterministic oblivious sequences.
All random choices defining the sequence are sampled before the game
starts.  We prove that the regret averaged over this distribution is at
least the desired quantity; hence at least one deterministic sequence in
its support has at least that much regret.

For $t=2,\ldots,T$, set
\[
\varepsilon_t:=\frac{1}{2t}.
\]
Let $Z_2,\ldots,Z_T$ be independent Bernoulli random variables satisfying
\[
\Pr(Z_t=1)=\varepsilon_t.
\]
Define the random price sequence $(P_t)_{t=1}^T$ by
\[
P_1:=1,
\qquad
P_t:=\prod_{s=2}^t(1-\varepsilon_s)^{Z_s},
\qquad t\ge2.
\]
Equivalently, conditional on $P_{t-1}=x$,
\[
P_t=
\begin{cases}
 x, & \text{with probability }1-\varepsilon_t,\\[1mm]
 (1-\varepsilon_t)x, & \text{with probability }\varepsilon_t.
\end{cases}
\tag{1}
\label{eq:transition}
\]
Every realization is nonincreasing.  Therefore, for every realization
$\bp_T$,
\[
M_t=p_t
\qquad\text{for every }t,
\]
and consequently
\[
\cV_T=0.
\tag{2}
\label{eq:variation-zero}
\]

\begin{lemma}
\label{lem:one-round}
Fix $t\ge2$.  Conditional on any price history with $P_{t-1}=x$, the
conditional expected reward of the learner at round $t$ is at most
\[
(1-\varepsilon_t)x.
\]
\end{lemma}

\begin{proof}
Let $A_t=a$.
There are three cases.

If $a<(1-\varepsilon_t)x$, then the order executes in both possible
states, and its reward is
\[
a\le(1-\varepsilon_t)x.
\]
If
\[
(1-\varepsilon_t)x\le a<x,
\]
then the order executes only when $P_t = P_{t-1}$, and its conditional
expected reward is
\[
(1-\varepsilon_t)a
\le
(1-\varepsilon_t)x.
\]
Finally, if $a\ge x$, the order never executes and its reward is zero.
\end{proof}

Let
\[
m_t:=\E[P_t],
\qquad t=1,\ldots,T,
\]
and introduce the convenient conventions
\[
m_0:=1,
\qquad
\varepsilon_1:=0.
\]

\begin{lemma}
\label{lem:price-path}
For every $t=1,\ldots,T$,
\[
m_t=(1-\varepsilon_t^2)m_{t-1}.
\tag{3}
\label{eq:mean-recursion}
\]
Moreover,
\[
m_t\ge\frac34
\qquad\text{for every }t=0,\ldots,T.
\tag{4}
\label{eq:mean-lower-bound}
\]
\end{lemma}

\begin{proof}
The claim is immediate for $t=1$.  For $t\ge2$, conditioning on
$P_{t-1}=x$ and using \eqref{eq:transition} gives
\[
\begin{aligned}
\E[P_t\mid P_{t-1}=x]
&=(1-\varepsilon_t)x
  +\varepsilon_t(1-\varepsilon_t)x\\
&=(1-\varepsilon_t^2)x.
\end{aligned}
\]
Taking expectations proves \eqref{eq:mean-recursion}.

Iterating the recursion gives
\[
m_t
=
\prod_{s=2}^t\left(1-\frac{1}{4s^2}\right).
\]
For numbers $x_i\in[0,1]$, one has
$\prod_i(1-x_i)\ge1-\sum_i x_i$.  Hence
\[
\begin{aligned}
m_t
&\ge
1-\sum_{s=2}^t\frac{1}{4s^2}\\
&\ge
1-\frac14\sum_{s=2}^{\infty}\frac{1}{s(s-1)}\\
&=1-\frac14
=
\frac34,
\end{aligned}
\]
which proves \eqref{eq:mean-lower-bound}.
\end{proof}

\begin{lemma}
\label{lem:benchmark}
For every realization $\bp_T$ of the random sequence,
\[
B_T(\bp_T)\ge T p_T.
\]
\end{lemma}

\begin{proof}
Because $\bp_T$ is nonincreasing, for every $a<p_T$ one has
$a<M_t=p_t$ for every $t$.  Therefore
\[
\sum_{t=1}^T r_t(a)=Ta.
\]
Taking the supremum over $a<p_T$ yields
\[
B_T(\bp_T)\ge\sup_{a<p_T}Ta=Tp_T.
\]
\end{proof}

We now combine the three lemmas.  At round $1$, the learner's reward is
at most $1=m_0$.  For every $t\ge2$, Lemma~\ref{lem:one-round} gives
\[
\E\bigl[r_t(A_t)\bigr]
\le
(1-\varepsilon_t)m_{t-1}.
\]
Thus, with the convention $\varepsilon_1=0$, the learner's expected
cumulative reward, where the expectation is now over both the random
sequence and the learner's internal randomization, is at most
\[
\sum_{t=1}^T(1-\varepsilon_t)m_{t-1}.
\tag{5}
\label{eq:learner-upper-bound}
\]

By Lemma~\ref{lem:benchmark}, the expected benchmark is at least $Tm_T$.
Consequently, the regret averaged over the random price sequence satisfies
\[
\E[R_T]
\ge
Tm_T-
\sum_{t=1}^T(1-\varepsilon_t)m_{t-1}.
\tag{6}
\label{eq:regret-start}
\]

We use the telescoping identity
\[
\sum_{t=1}^T m_{t-1}-Tm_T
=
\sum_{t=1}^T t\bigl(m_{t-1}-m_t\bigr).
\tag{7}
\label{eq:telescoping}
\]
Indeed, expanding the right-hand side gives
$m_0+m_1+\cdots+m_{T-1}-Tm_T$.
By \eqref{eq:mean-recursion},
\[
m_{t-1}-m_t
=
\varepsilon_t^2m_{t-1}.
\tag{8}
\label{eq:mean-drop}
\]
Substituting \eqref{eq:telescoping} and \eqref{eq:mean-drop} into
\eqref{eq:regret-start} yields
\[
\begin{aligned}
\E[R_T]
&\ge
\sum_{t=1}^T
\bigl(\varepsilon_t-t\varepsilon_t^2\bigr)m_{t-1}\\
&=
\sum_{t=2}^T
\left(\frac{1}{2t}-\frac{1}{4t}\right)m_{t-1}\\
&=
\frac14\sum_{t=2}^T\frac{m_{t-1}}{t}.
\end{aligned}
\tag{9}
\label{eq:harmonic-sum}
\]
Using Lemma~\ref{lem:price-path},
\[
\E[R_T]
\ge
\frac14\cdot\frac34\sum_{t=2}^T\frac1t
=
\frac{3}{16}\bigl(H_T-1\bigr).
\tag{10}
\label{eq:average-lower-bound}
\]

Therefore at least one deterministic sequence
$\bp_T$ in the support satisfies
\[
R_T
\ge
\frac{3}{16}\bigl(H_T-1\bigr).
\]
This proves the theorem.
\end{proof}

\end{document}
